% TEMPLATE-MILC-v3.tex - plantilla maestra para todas las guias MILC v3
% Ver scripts/generadores/build-guia-milc.py para la lista completa de
% claves esperadas. El script valida que no queden placeholders sin
% reemplazar antes de escribir el archivo final.

\documentclass[11pt,letterpaper]{article}
\usepackage[margin=1.75cm]{geometry}
\usepackage[table]{xcolor}
\usepackage{fontspec}
\usepackage[spanish]{babel}
\usepackage{tikz}
\usepackage[most]{tcolorbox}
\usepackage{tabularx}
\usepackage{array}
\usepackage{fancyhdr}
\usepackage{titlesec}
\usepackage{ragged2e}
\usepackage{enumitem}
\usepackage{microtype}

\setmainfont{Helvetica Neue}
\setsansfont{Helvetica Neue}
\setlength{\parindent}{0pt}
\setlength{\parskip}{6pt}
\hyphenpenalty=200
\exhyphenpenalty=200
\pretolerance=400
\tolerance=2000
\setlength{\emergencystretch}{1em}
\renewcommand{\arraystretch}{1.25}
\setlist{leftmargin=5mm,itemsep=1mm,topsep=1mm}
\newcolumntype{Y}{>{\RaggedRight\arraybackslash}X}

\definecolor{milcMagenta}{HTML}{E6007E}
\definecolor{milcVino}{HTML}{5A0038}
\definecolor{milcTurquesa}{HTML}{0093A5}
\definecolor{milcVerde}{HTML}{58A923}
\definecolor{milcMostaza}{HTML}{E5B400}
\definecolor{milcOcre}{HTML}{B86600}
\definecolor{milcNegro}{HTML}{241816}
\definecolor{milcGris}{HTML}{F7F7F4}
\definecolor{milcLinea}{HTML}{E8E2DD}
\definecolor{milcGradePrimary}{HTML}{7C3AED}
\definecolor{milcGradeDark}{HTML}{4C1D95}
\definecolor{milcGradeSoft}{HTML}{EDE9FE}
\definecolor{milcGradeAccent}{HTML}{CFE7FF}
\definecolor{milcGradeText}{HTML}{FFFFFF}
\color{milcNegro}

\pagestyle{fancy}
\fancyhf{}
\lhead{\small Tecnología e Informática · Grado 9}
\rhead{\small Guía 14 · MILC v3}
\cfoot{\small\thepage}
\renewcommand{\headrulewidth}{0pt}

\newtcolorbox{titlebox}[1]{
  enhanced,colback=#1,colframe=#1,arc=7mm,boxrule=0pt,
  left=7mm,right=7mm,top=3mm,bottom=3mm,
  before skip=8pt,after skip=8pt,
  fontupper=\sffamily\bfseries\Large\color{white}
}

\newtcolorbox{softbox}[3]{
  enhanced,colback=#2,colframe=#1,borderline west={4pt}{0pt}{#1},
  arc=2mm,boxrule=.4pt,left=4mm,right=4mm,top=3mm,bottom=3mm,
  before skip=6pt,after skip=8pt,title={#3},coltitle=white,
  colbacktitle=#1,fonttitle=\sffamily\bfseries,fontupper=\RaggedRight,
  attach boxed title to top left={xshift=3mm,yshift=-2mm},
  boxed title style={arc=2mm, boxrule=0pt}
}

\newtcolorbox{drawbox}[2]{
  enhanced,colback=white,colframe=#1,arc=4mm,boxrule=1pt,
  left=5mm,right=5mm,top=4mm,bottom=4mm,before skip=8pt,after skip=8pt,
  title={#2},coltitle=white,colbacktitle=#1,fonttitle=\sffamily\bfseries,
  fontupper=\RaggedRight,
  attach boxed title to top left={xshift=4mm,yshift=-2mm},
  boxed title style={arc=3mm, boxrule=0pt}
}

\newtcolorbox{infoband}[1]{
  enhanced,colback=#1!8,colframe=#1!50,arc=2mm,boxrule=.5pt,
  left=4mm,right=4mm,top=2mm,bottom=2mm,before skip=4pt,after skip=4pt,
  fontupper=\small\RaggedRight
}

% Puente narrativo entre secciones (contrato editorial Conéctate).
% Da continuidad: "Acabas de X. Ahora vas a Y porque Z."
% Visualmente: banda izquierda en color del grado, texto en itálica pequeña.
\newtcolorbox{puentebox}{
  enhanced,colback=milcGradeSoft,colframe=milcGradeDark,
  borderline west={3pt}{0pt}{milcGradeDark},
  arc=0mm,boxrule=0pt,left=5mm,right=4mm,top=2.5mm,bottom=2.5mm,
  before skip=4pt,after skip=8pt,
  fontupper=\itshape\small\color{milcNegro}\RaggedRight
}
\newcommand{\puente}[1]{%
  \begin{puentebox}%
    \textcolor{milcGradeDark}{\textbf{\textrightarrow}}\hspace{2mm}#1%
  \end{puentebox}%
}

\newcommand{\linead}[1]{\noindent\color{milcLinea}\rule{#1}{0.5pt}\color{milcNegro}}
\newcommand{\respuesta}[1]{\par\vspace{2mm}\foreach \n in {1,...,#1}{\noindent\color{milcLinea}\rule{\linewidth}{0.5pt}\par\vspace{5mm}}\color{milcNegro}}
\newcommand{\checkbox}{\textcolor{milcGradeDark}{\fbox{\phantom{x}}}\hspace{5pt}}

\newcommand{\phasecard}[4]{
\begin{tcolorbox}[enhanced,colback=#3,colframe=#2,arc=3mm,boxrule=.5pt,height=3.05cm,valign=center,halign=center,left=1.5mm,right=1.5mm,top=2mm,bottom=2mm]
{\sffamily\bfseries\fontsize{22}{24}\selectfont\textcolor{#2}{#1}}\par
{\sffamily\bfseries\small #4}
\end{tcolorbox}
}

\begin{document}
\thispagestyle{empty}
\begin{tikzpicture}[remember picture,overlay]
  \fill[white] (current page.south west) rectangle (current page.north east);
  \path[fill=milcGradePrimary,rounded corners=16mm]
    ([xshift=.55cm,yshift=-.55cm]current page.north west) rectangle
    ([xshift=-.55cm,yshift=3.65cm]current page.south east);

  \node[anchor=north west,text=milcGradeText] at ([xshift=2.0cm,yshift=-1.85cm]current page.north west)
    {\sffamily\bfseries\fontsize{14}{16}\selectfont GRADO};
  \node[anchor=north west,text=milcGradeText,opacity=.82] at ([xshift=2.0cm,yshift=-2.75cm]current page.north west)
    {\sffamily\bfseries\fontsize{9}{11}\selectfont SERIE GUÍAS · TECNOLOGÍA E INFORMÁTICA};

  \begin{scope}[shift={([xshift=-3.05cm,yshift=-2.55cm]current page.north east)}]
    \fill[white,opacity=.80] (-.62,-.52) rectangle (.62,.52);
    \draw[milcGradeText,line width=1pt] (-.62,-.52) rectangle (.62,.52);
    \fill[milcGradeDark] (-.42,-.38) rectangle (-.22,.06);
    \fill[milcGradeAccent] (-.10,-.38) rectangle (.10,.24);
    \fill[milcGradePrimary!60!black] (.22,-.38) rectangle (.42,.38);
  \end{scope}

  \node[anchor=west,text=milcGradeText] at ([xshift=1.9cm,yshift=-12.85cm]current page.north west)
    {\sffamily\bfseries\fontsize{124}{124}\selectfont 9};
  \draw[milcGradeText,line width=1.7pt]
    ([xshift=4.52cm,yshift=-11.68cm]current page.north west) circle (.13cm);

  \node[anchor=west,text=milcGradeText,text width=8.7cm,align=left] at ([xshift=10.35cm,yshift=-11.6cm]current page.north west)
    {
     {\sffamily\bfseries\fontsize{19}{22}\selectfont Color en editorial ---\\contraste, jerarquía,\\accesibilidad}\\[4mm]
     {\sffamily\bfseries\fontsize{23}{26}\selectfont Guía 14-9-TIC}\\[2mm]
     {\sffamily\fontsize{13}{16}\selectfont Período 2 · Diseño editorial digital}\\[5mm]
     {\sffamily\bfseries\fontsize{11}{13}\selectfont Institución Educativa Sor María Juliana}\\[1mm]
     {\sffamily\fontsize{10}{12}\selectfont Autor: Dr. Álvaro Cárdenas Orozco}
    };

  \path[fill=milcGradeSoft,rounded corners=7mm]
    ([xshift=1.35cm,yshift=.85cm]current page.south west) rectangle
    ([xshift=-1.35cm,yshift=4.25cm]current page.south east);
  \draw[milcGradeDark,line width=.65pt,rounded corners=7mm]
    ([xshift=1.35cm,yshift=.85cm]current page.south west) rectangle
    ([xshift=-1.35cm,yshift=4.25cm]current page.south east);
  \node[anchor=center,text=milcNegro,text width=17.0cm,align=center] at ([yshift=2.55cm]current page.south)
    {\sffamily\fontsize{10}{12}\selectfont Metodología MILC · Apertura ancestral · Pensamiento computacional · Triángulo Dussel-Estoicismo-Floridi\\
    \textcolor{milcGradeDark}{\bfseries Producto:} Paleta editorial de 5 colores con código HEX y función declarada (dominante, secundario, acento, fondo, texto), justificada en 3 frases sobre la emoción que transmite a la revista};
\end{tikzpicture}
\mbox{}
\newpage

\section*{Apertura: Color en editorial --- contraste, jerarquía, accesibilidad}

\begin{softbox}{milcGradeDark}{milcGradeSoft}{Saber ancestral}
Las paletas naturales del Valle, el Pacífico y los Andes son ingeniería estética milenaria. El \textbf{verde del café} en la zona cafetera; el \textbf{ocre del sombrero vueltiao}; el \textbf{rojo del achiote} con que se pintaban guerreros y se pintan hoy alimentos; el \textbf{azul del Pacífico} en el horizonte de Buenaventura y Tumaco; el \textbf{negro brillante de las molas guna} con sus contrastes precisos. Cada paleta tradicional fue probada por siglos de uso humano: armoniosa, simbólica, reconocible. La paleta no nació en Pantone --- nació donde la tierra, las plantas y el cielo se miraron juntos.
\end{softbox}

\begin{softbox}{milcTurquesa}{milcGris}{Saber contemporáneo}
Hoy una \textbf{paleta editorial} son 3-5 colores con códigos HEX y \emph{función declarada}: uno dominante, uno secundario, uno o dos acentos, uno de fondo, uno de texto. Cada color cumple un trabajo. Si todos son ``bonitos'' pero ninguno tiene función clara, la paleta se vuelve confusión. Diseñar con paletas locales (Valle, Pacífico) no es nostalgia: es \emph{coherencia visual con el lugar donde publicas}.
\end{softbox}

\begin{softbox}{milcMostaza}{milcGris}{Pregunta-puente}
\textit{Mira el sombrero vueltiao o una mola guna: ¿qué hace que esos colores combinen perfecto aunque sean intensos? ¿Y qué falla en muchas paletas digitales que parecen ``saturadas'' o ``confusas''?}
\end{softbox}

\begin{softbox}{milcVerde}{milcGris}{Saber y saber hacer}
Al terminar podrás: (1) \textbf{identificar} 5 paletas reales (naturales, tradicionales, comerciales); (2) \textbf{analizar} los 3 principios del color editorial: contraste, jerarquía, accesibilidad; (3) \textbf{aplicar} esos principios a una paleta de 5 colores; (4) \textbf{crear} la paleta editorial de tu revista con función declarada y justificación.
\end{softbox}

\small
\begin{tabularx}{\textwidth}{>{\bfseries}p{3.0cm}Y}
\rowcolor{milcGradePrimary!12}Autor & Dr. Álvaro Cárdenas Orozco\\
DBA & Producción de piezas editoriales digitales con criterio estético y ético (MEN, T\&I)\\
Referentes & Tipografía colombiana · Diseño centrado en accesibilidad · Floridi\\
Producto final & Paleta editorial de 5 colores con código HEX y función declarada (dominante, secundario, acento, fondo, texto), justificada en 3 frases sobre la emoción que transmite a la revista\\
\end{tabularx}
\normalsize

\puente{Antes de empezar, mira el viaje completo. La sesión 3 te dio tipografía; hoy le agregas color a tu revista. Con grilla + tipografía + color tienes el \emph{sistema visual mínimo} para diseñar cualquier pieza editorial.}

\begin{titlebox}{milcGradeDark}
Ruta MILC: así vamos a aprender
\end{titlebox}
\begin{tabularx}{\textwidth}{>{\centering\arraybackslash}X>{\centering\arraybackslash}X>{\centering\arraybackslash}X>{\centering\arraybackslash}X}
\phasecard{1}{milcVerde}{milcGris}{Escuta\\[-1mm]\small Observo, escucho y pregunto.} &
\phasecard{2}{milcTurquesa}{milcGris}{Sistematización\\[-1mm]\small Pensamiento computacional.} &
\phasecard{3}{milcMagenta}{milcGris}{Praxis\\[-1mm]\small Construyo y aplico.} &
\phasecard{4}{milcVino}{milcGris}{Evaluación\\[-1mm]\small Reflexión liberadora.}
\end{tabularx}


\puente{Antes de elegir tu paleta, vamos a leer paletas reales que llevan siglos funcionando: la del Valle cafetero, la del Pacífico, la wayuu, la de marcas globales. Cada una tiene lógica detrás.}

\begin{titlebox}{milcVerde}
Fase 1 · Escuta
\end{titlebox}

\begin{softbox}{milcVerde}{milcGris}{Escena para escuchar}
\textbf{Actividad 1 · IDENTIFICA --- 5 paletas reales} (15 min · individual). Busca \textbf{5 paletas distintas}: (a) una \emph{natural} (paisaje de tu región: café, montaña, Pacífico, llanos); (b) una \emph{tradicional} (sombrero vueltiao, mola guna, manta wayuu, sayo del Pacífico); (c) una de \emph{marca global} (Coca-Cola, Spotify, Apple); (d) una de \emph{red social} que usas; (e) una \emph{paleta literaria} (la portada de un libro que te guste). Para cada una: ¿cuántos colores predomina? ¿qué emoción transmite?
\end{softbox}

\checkbox Identifiqué 5 paletas distintas (natural + tradicional + global + redes + literaria).\\
\checkbox Conté colores principales en cada una.\\
\checkbox Anoté la emoción que transmite cada paleta.

\begin{infoband}{milcVerde}
\textbf{Lo que va al cuaderno (Actividad 1 · IDENTIFICA).} Título: «Actividad 1 --- 5 paletas que comunican». \textbf{Formato:} tabla de 4 columnas (Paleta~|~Colores predominantes~|~Cantidad aprox.~|~Emoción transmitida), 5 filas. \textbf{Sabes que terminaste cuando} reconoces patrones (las paletas fuertes casi nunca usan más de 5 colores; los neón son raros).
\end{infoband}

\puente{Ya viste paletas reales. Ahora vamos a entender los \textbf{3 principios del color editorial}: contraste (para distinguir), jerarquía (para ordenar la lectura), accesibilidad (para que todos puedan leer). Sin alguno, la paleta puede verse bonita pero falla en función.}

\begin{titlebox}{milcTurquesa}
Fase 2 · Sistematización con pensamiento computacional
\end{titlebox}

\begin{softbox}{milcTurquesa}{milcGris}{Cuatro pilares aplicados al tema}
Toda paleta editorial fuerte cumple 3 principios. Vamos a descomponerlos con los pilares del pensamiento computacional para que tu elección sea defendible y no por gusto.
\end{softbox}

\small
\begin{tabularx}{\textwidth}{>{\bfseries\RaggedRight\arraybackslash}p{3.6cm}Y}
\rowcolor{milcTurquesa!90}\color{white}Pilar & \color{white}\bfseries Aplicación al tema\\
1. Descomposición & El color editorial se descompone en 3 principios: (1) \textbf{contraste} (los colores deben distinguirse claramente entre sí --- contraste tonal, no solo cromático), (2) \textbf{jerarquía} (un color dominante 60\%, uno secundario 30\%, acentos 10\% para guiar el ojo), (3) \textbf{accesibilidad} (combinaciones que cumplen contraste mínimo WCAG AA para que personas con baja visión o daltonismo puedan leer).\\
2. Reconocimiento de patrones & Toda paleta fuerte sigue el patrón \textbf{``regla 60-30-10''}: dominante 60\% (color principal), secundario 30\% (apoyo, fondos), acento 10\% (botones, llamadas a la acción). Cambia las proporciones y la paleta pierde armonía. El acento que ocupa 50\% deja de ser acento.\\
3. Abstracción & Lo esencial cabe en una pregunta: \emph{¿esta combinación se lee bien para alguien con baja visión, daltonismo o pantalla mala?}. Si la respuesta es no, falla aunque sea ``bonita''. El diseño verdaderamente bonito incluye.\\
4. Algoritmo conceptual & Pasos para construir una paleta: (1) define emoción central de tu revista; (2) elige UN dominante; (3) elige UN secundario complementario o análogo; (4) elige 1-2 acentos contrastantes; (5) verifica con webaim.org/resources/contrastchecker que el contraste texto-fondo sea $\geq$ 4.5:1.\\
\end{tabularx}
\normalsize

\begin{drawbox}{milcTurquesa}{Anatomía de una paleta editorial}
\textbf{Color dominante (60\%):} el color principal de la marca o revista. \\[1mm] \textbf{Color secundario (30\%):} apoyo del dominante, para fondos secundarios o cuerpos de texto. \\[1mm] \textbf{Acento (10\%):} botones, llamadas a la acción, destacados. Suele contrastar. \\[1mm] \textbf{Fondo:} casi siempre blanco o muy claro (legibilidad). \\[1mm] \textbf{Texto:} negro o muy oscuro sobre fondo claro. \\[1mm] \textbf{Códigos HEX:} todos anotados con \#HHHHHH para reproducirlos.
\end{drawbox}

\begin{softbox}{milcMostaza}{milcGris}{Errores comunes y su corrección}
(1) Demasiados colores (>5): paleta confusa. (2) Todos con misma saturación: sin jerarquía visual. (3) Texto gris claro sobre fondo blanco: ilegible para muchas personas. (4) Imitar paleta de marca famosa sin entender por qué funciona. (5) Olvidar que el daltonismo afecta al 8\% de hombres --- combinaciones rojo/verde fallan para ellos.
\end{softbox}

\begin{infoband}{milcTurquesa}
\textbf{Lo que va al cuaderno (Actividad 2 · EVALÚA).} Título: «Actividad 2 --- Diagnóstico de paletas reales». \textbf{Formato:} 3 fichas, una por principio (contraste, jerarquía, accesibilidad), cada una con: definición en tus palabras + ejemplo de una de las 5 paletas de la Actividad 1 que lo cumple bien + ejemplo de algo que lo incumple. \textbf{Veredicto:} mira la portada de un libro o revista que tengas a la mano. ¿Aprueba o reprueba los 3 principios? Si reprueba alguno, ¿cuál es el principio incumplido más grave y por qué? Justifica con criterios objetivos (no gusto). \textbf{Sabes que terminaste cuando} puedes diagnosticar problemas de color en cualquier diseño con argumentos verificables.
\end{infoband}

\puente{Tienes el método. Toca elegir: tu paleta editorial de 5 colores para la revista del periodo, con HEX y función declarada para cada uno.}

\begin{titlebox}{milcMagenta}
Fase 3 · Praxis con pensamiento computacional
\end{titlebox}

\begin{softbox}{milcMagenta}{milcGris}{Construyendo el producto con los cuatro pilares}
\textbf{Actividad 3 · CREA --- Tu paleta editorial} (30 min · individual). Hora de elegir. Vas a construir la paleta de 5 colores para tu revista del periodo, con códigos HEX, función declarada para cada uno, y justificación de 3 frases sobre la emoción que transmite.
\end{softbox}

\small
\begin{tabularx}{\textwidth}{>{\bfseries\RaggedRight\arraybackslash}p{3.6cm}Y}
\rowcolor{milcMagenta!90}\color{white}Pilar & \color{white}\bfseries Aplicación al producto\\
Descomposición & Tu paleta se descompone en 5 colores con función: (1) dominante (60\%, color principal), (2) secundario (30\%, apoyo), (3) acento 1 (10\%, llamadas a la acción), (4) fondo (casi siempre blanco o muy claro), (5) texto (casi siempre negro o muy oscuro). Cada uno con su código HEX.\\
Patrones & Sigue el patrón \textbf{``natural antes que tendencia''}: tu primera versión puede inspirarse en una paleta natural local (café, mango, achiote, mar Pacífico, cielo del Valle). Eso le da identidad inmediata y evita la trampa de copiar marcas globales.\\
Abstracción & Cada color expresa una sola función. No hagas que tu acento sirva también de fondo --- elige uno para cada cosa.\\
Algoritmo de armado & Algoritmo: (1) declara emoción central de tu revista; (2) usa coolors.co o pinta a mano 3 paletas candidatas; (3) elige UNA; (4) anota HEX de los 5 colores; (5) declara función para cada uno; (6) verifica contraste texto-fondo en webaim.org; (7) escribe 3 frases de justificación.\\
\end{tabularx}
\normalsize

\begin{drawbox}{milcMagenta}{Checklist antes de cerrar la paleta}
\checkbox 5 colores con código HEX anotado. \\[1mm] \checkbox Función declarada para cada uno (dominante / secundario / acento / fondo / texto). \\[1mm] \checkbox Contraste texto-fondo verificado en webaim.org ($\geq$ 4.5:1). \\[1mm] \checkbox Al menos UN color inspirado en paleta natural o tradicional local. \\[1mm] \checkbox Justificación de 3 frases sobre emoción transmitida. \\[1mm] \checkbox La paleta se siente coherente al pasar la mirada por todos los colores juntos.
\end{drawbox}

\begin{softbox}{milcMagenta}{milcGris}{Plantilla de trabajo}
\textbf{Emoción central de mi revista.} \textit{1 palabra.} \\[1mm] \textbf{Mi paleta.} \textit{5 colores con HEX y función.} \\[1mm] \textbf{Inspiración local.} \textit{Qué paleta natural o tradicional inspira uno o más colores.} \\[1mm] \textbf{Justificación.} \textit{3 frases sobre por qué esa paleta sirve a mi revista.}
\end{softbox}

\begin{infoband}{milcMagenta}
\textbf{Lo que va al cuaderno (Actividad 3 · CREA).} Título: «Actividad 3 --- Mi paleta editorial». \textbf{Formato:} 5 cuadritos pintados o pegados con HEX debajo + función declarada + justificación de 3 frases. \textbf{Extensión:} 1 página de cuaderno. \textbf{Sabes que terminaste cuando} puedes defender tu paleta frente a un compañero sin decir ``porque me gusta''.
\end{infoband}

\puente{Lo que sigue es el resumen de qué entregas hoy: paleta de 5 colores + justificación. El criterio principal: que cada color tenga función clara y la paleta pase el test de accesibilidad (contraste mínimo).}

\begin{titlebox}{milcMostaza}
Producto de la sesión
\end{titlebox}
\textbf{Paleta editorial de 5 colores con HEX, función declarada y justificación}.
\begin{softbox}{milcMostaza}{milcGris}{Criterios de calidad}
(1) 5 colores con código HEX anotado. (2) Función declarada para cada uno. (3) Contraste texto-fondo verificado en webaim.org ($\geq$ 4.5:1). (4) Al menos 1 color inspirado en paleta natural o tradicional local. (5) Justificación de 3 frases sobre emoción transmitida. (6) Pasa la prueba del "test del extraño" (sin contexto, transmite la emoción declarada).
\end{softbox}

\puente{Antes de cerrar, mira la elección de color desde las cinco dimensiones humanas. El color no es solo estética: es ciudadanía digital (¿quién puede leer mi sitio?) y soberanía cultural (¿qué paleta heredo o invento?).}

\begin{titlebox}{milcVino}
Fase 4 · Evaluación liberadora — cinco dimensiones
\end{titlebox}

\begin{softbox}{milcVino}{milcGris}{Sentido de la evaluación}
Evaluar no es solo revisar una nota. Es reconocer qué aprendí, qué cambió en mí, qué emociones manejé, cómo me relaciono con la ciudad y la comunidad, y qué le dejaré a quien viene detrás.
\end{softbox}

\textbf{Mi reflexión liberadora en cinco dimensiones.} Completa cada frase iniciada:

\small
\begin{tabularx}{\textwidth}{>{\bfseries\RaggedRight\arraybackslash}p{3.4cm}Y>{\RaggedRight\arraybackslash}p{4.6cm}}
\rowcolor{milcVino}\color{white}Dimensión & \color{white}\bfseries Mi respuesta (frame starter) & \color{white}\bfseries Algo que descubrí\\
1. Personal & Lo que más cambió en mí fue \linead{6.5cm} & \linead{4.2cm}\\
2. Emocional & La emoción más fuerte fue \linead{2.5cm} y la regulé \linead{3cm} & \linead{4.2cm}\\
3. Ciudadana & En democracia esto importa porque \linead{6cm} & \linead{4.2cm}\\
4. Local (Cartago) & En mi barrio sería útil para \linead{6.5cm} & \linead{4.2cm}\\
5. Intergeneracional & A mi abuela/abuelo le diría \linead{6.5cm} & \linead{4.2cm}\\
\end{tabularx}
\normalsize

\begin{infoband}{milcVino}
\textbf{Para la próxima sustentación.} Vuelve a esta tabla en seis meses. Verás que las respuestas cambian — no porque lo aprendido cambie, sino porque tú cambias. La evaluación liberadora es una conversación contigo a través del tiempo.
\end{infoband}


\puente{Y para terminar, tres voces sobre color. Dussel sobre las paletas colonizadas, Marco Aurelio sobre la sobriedad como elegancia, Floridi sobre el color como infraestructura de la infoesfera.}

\begin{titlebox}{milcGradeDark}
Triángulo de pensamiento · cierre filosófico
\end{titlebox}

\begin{softbox}{milcVino}{milcGris}{Dussel · lente del nosotros}
\textit{``Las paletas que dominan la infoesfera global no son universales --- son culturalmente específicas y se imponen.''}

El blanco minimalista de las marcas tech (Apple, Notion, Linear) viene de la modernidad escandinava del siglo XX. La paleta tropical (verde palmera + azul caribe + naranja sunset) viene del marketing turístico. Cuando eliges paleta sin pensar, eliges una cultura. Atreverte a una paleta que mezcle achiote, ocre y verde café es declarar que tu revista habla desde tu lugar, no desde un genérico global.

\textbf{¿Mi paleta copia un estándar global (tech minimalista, tropical turístico) o se atreve a una identidad local? ¿Qué precio paga cada elección?}
\end{softbox}
\linead{\linewidth}\\[1mm]
\linead{\linewidth}

\begin{softbox}{milcMostaza}{milcGris}{Estoicismo · Marco Aurelio · cuidado interior}
\textit{``Pocos colores bien elegidos comunican más que muchos colores mal organizados.''}

La tentación del diseñador novato: usar más colores ``para que se vea vivo''. La sabiduría estoica del diseño: 3-5 colores con función clara comunican más que 12 colores sin jerarquía. The Economist usa básicamente 3 colores (rojo, blanco, negro) y lleva 180 años. Apple usa 2 (blanco y gris). Lo sobrio dura.

\textbf{¿Por qué siento la tentación de agregar más colores cuando algo se ve ``aburrido''? ¿Será que lo aburrido es sobrio bien hecho?}
\end{softbox}
\linead{\linewidth}\\[1mm]
\linead{\linewidth}

\begin{softbox}{milcTurquesa}{milcGris}{Floridi · lente de la infoesfera}
\textit{``El color es la infraestructura cromática de la infoesfera --- decide quién puede leer y quién queda afuera.''}

Una combinación de colores que tú ves como ``elegante'' puede ser invisible para alguien con daltonismo. Una página gris claro sobre blanco puede ser legible en tu pantalla de 1500\$ pero ilegible en una pantalla de gama básica. Diseñar con accesibilidad es ciudadanía digital: incluyes o excluyes con cada decisión cromática.

\textbf{¿Mi paleta funciona para alguien con daltonismo? ¿Para alguien que mira mi revista en un celular básico bajo el sol? ¿A quién dejo afuera con mi elección?}
\end{softbox}
\linead{\linewidth}\\[1mm]
\linead{\linewidth}

\begin{titlebox}{milcVino}
Compromiso final
\end{titlebox}

\small
\begin{tabularx}{\textwidth}{>{\RaggedRight\arraybackslash}p{3.0cm}YYY}
\rowcolor{milcVino}\color{white}\bfseries Criterio & \color{white}\bfseries Lo logré & \color{white}\bfseries En proceso & \color{white}\bfseries Necesito apoyo\\
Comprensión del tema & Explico ideas centrales con mis palabras. & Reconozco ideas pero falta precisión. & Necesito repasar conceptos base.\\
Pensamiento computacional & Apliqué los 4 pilares al diseñar mi producto. & Apliqué algunos pilares. & Necesito releer la fase 2.\\
Aplicación y producto & Mi producto cumple los criterios y se puede revisar. & Producto funcional con ajustes pendientes. & Aún en construcción.\\
Reflexión 5 dimensiones & Respondí cada dimensión con honestidad. & Respondí algunas con detalle. & Mis respuestas son muy generales.\\
Triángulo de pensamiento & Conecté las 3 voces con mi trabajo real. & Conecté al menos una. & No relaciono las voces con mi proceso.\\
\end{tabularx}
\normalsize

\textbf{Hoy entendí que:}\\[1mm]
\linead{\linewidth}\\[1mm]
\linead{\linewidth}

\textbf{Mi compromiso para la próxima sesión es:}\\[1mm]
\linead{\linewidth}\\[1mm]
\linead{\linewidth}

\begin{softbox}{milcMostaza}{milcGris}{Cita de despedida · Marco Aurelio}
\textit{``El obstáculo en el camino se vuelve el camino. Lo que estorba al camino se vuelve camino.''}\\[1mm]
Cada error de hoy, bien observado, se convierte en pieza del camino que pisarás mañana.
\end{softbox}

\small
\begin{tabularx}{\textwidth}{>{\bfseries}p{3.6cm}Y}
\rowcolor{milcGradeDark!12}Nombre completo: & \linead{10cm}\\
Fecha: & \linead{4cm} \quad Curso/grupo: \linead{4cm}\\
Mi firma: & \linead{9cm}\\
\end{tabularx}
\normalsize

\begin{infoband}{milcGradeDark}
\textbf{Para la próxima semana.} Aplica tu paleta a una pieza real esta semana: una hoja de tu cuaderno, una historia de Instagram, un grupo de Whatsapp, un avatar nuevo. Anota qué reacción tuvo (interna o externa). Si nadie comentó nada y a ti te incomodó, ajusta. La paleta no es ley --- es una hipótesis que validas usándola. El próximo lunes traemos esa bitácora.
\end{infoband}

\end{document}
